% Generated draft: requires mathematical and source review.
\documentclass[11pt]{article}
\usepackage[margin=25mm]{geometry}
\usepackage{fontspec}
\setmainfont{Latin Modern Roman}
\usepackage{amsmath,amssymb,mathtools,unicode-math}
\setmathfont{Latin Modern Math}
\usepackage{enumitem,needspace,xurl,hyperref}
\setlength{\parindent}{0pt}
\setlength{\parskip}{5pt}
\setlength{\emergencystretch}{3em}
\newcommand{\N}{\mathbb{N}}
\newcommand{\Z}{\mathbb{Z}}
\newcommand{\Q}{\mathbb{Q}}
\newcommand{\R}{\mathbb{R}}
\newcommand{\C}{\mathbb{C}}
\newcommand{\F}{\mathbb{F}}
\newcommand{\A}{\mathbb{A}}
\newcommand{\PP}{\mathbb{P}}
\newcommand{\E}{\mathbb{E}}
\newcommand{\eps}{\varepsilon}
\newcommand{\dd}{\,\mathrm{d}}
\begin{document}
\section*{UnsolvedMath ID: 800963}
\textbf{Source problem code:} catalogue entry
\subsection{Definitions and mathematical statement}
Let \((M,g)\) be an oriented closed Riemannian surface of genus at least \(1\), and let \(d_g\) be its Riemannian distance. A nontrivial free homotopy class \(\Gamma\) is a free homotopy class of loops in \(M\) that does not contain a constant loop; no primitivity assumption is stated. A closed geodesic written as \(\gamma:\R\to M\) is understood through its periodic extension: for some \(T>0\), \(\gamma(t+T)=\gamma(t)\), and the loop \(s\mapsto\gamma(Ts)\), with \(s\in\R/\Z\), represents \(\Gamma\).

The supplied source attributes the following surface assertion to Gaidukov: for every such \((M,g)\), every nontrivial free homotopy class \(\Gamma\), and every \(p\in M\), there exist a closed geodesic \(\gamma:\R\to M\) representing \(\Gamma\) and a ray \(\beta:[0,\infty)\to M\) such that
\[
\beta(0)=p,
\qquad
\lim_{t\to\infty}d_g\bigl(\beta(t),\gamma(t)\bigr)=0.
\]
This is convergence at matching parameter values, rather than merely convergence of the distance to the image of \(\gamma\).

\textbf{Requested generalization.} Generalize this surface assertion to higher dimensions.

\textbf{Unresolved formulation qualifications.} The supplied record does not specify the higher-dimensional manifold class, additional geometric or topological hypotheses, or whether the conclusion should remain exactly the displayed convergence assertion. Thus it does not determine a single precise higher-dimensional conjecture. It also does not explicitly define its convention for a ray or specify the speed normalization and parametrization of the two curves. If a geodesic ray is intended, global distance minimization in \(M\) cannot be required for an infinite unit-speed ray on a closed manifold, since its diameter is finite. The intended convention must be checked in the cited source before a fully specified generalization is formulated.
\subsection{Short English statement}
The supplied surface theorem says that, from every point of an oriented closed Riemannian surface of positive genus, a ray approaches a periodically parametrized closed geodesic in any prescribed nontrivial free homotopy class at matching times. The problem asks for a higher-dimensional generalization, but leaves its hypotheses and the precise ray convention unspecified.
\subsection{Sources}
\begin{itemize}
\item Supplied local UnsolvedMath registry record, permanent ID \(800963\), exact source code \texttt{catalogue entry}, in the AMR Open Problem Lists collection. The registry labels the problem open; that label is supplied metadata, not an independently verified mathematical or literature assessment.
\item Burns and Matveev, \emph{Open Problems and Questions about Geodesics} (2019), Problem 44, as identified by the supplied record: \url{https://arxiv.org/abs/1308.5417}. The record reports access on 2026-10-16 and extraction from arXiv LaTeX; these provenance details have not been independently verified here. The registry lists Burns and Matveev as proposers and gives Benjamin Schmidt as attribution/context.
\item Gaidukov, source reference label \texttt{[Gaidukov1966]}, cited in the supplied statement for the surface theorem. The same statement also spells the name Gajdukov. No full bibliographic entry is supplied.
\item Source reference label \texttt{[Bia-Pol1986]}. The supplied statement says this reference explains how Gaidukov's results follow easily from the classical results cited as \texttt{[Morse1924]} and \texttt{[Hedlund1932]}. Full bibliographic entries for these three references are not supplied.
\end{itemize}
\end{document}
